\chapter{Modelos de Contagem}
\label{cap:contagem}
% ── índice ──────────────────────────────────────────────────────
\index{modelos de contagem|textbf}
\index{Poisson, regressão de|textbf}
\index{poisson@\texttt{poisson()}|textbf}
\index{binomial negativa|textbf}
\index{nbreg@\texttt{nbreg()}|textbf}
\index{ZINB|textbf}
\index{zinb@\texttt{zinb()}|textbf}
\index{superdispersão|textbf}
\index{excesso de zeros}
\index{pseudo-máxima verossimilhança}

% ─────────────────────────────────────────────────────────────────────────────
\section{Dados de contagem e a inadequação do OLS}

Variáveis de resultado que assumem apenas valores inteiros não negativos
($Y \in \{0,1,2,\ldots\}$) — número de filhos, visitas ao médico, patentes
registradas — violam as premissas do OLS. O OLS pode produzir predições
negativas e é ineficiente: a variância de $Y$ cresce com a média (cauda
assimétrica à direita), violando homocedasticidade.

\subsection*{Poisson}

O modelo de regressão Poisson especifica:

\[
  Y_i \mid x_i \sim \mathrm{Poisson}(\lambda_i)
  \qquad
  \lambda_i = E[Y_i \mid x_i] = \exp(\beta_0 + \beta_1 x_i)
\]

A função de log-verossimilhança é:

\[
  \ell(\beta) = \sum_{i=1}^n \left[ y_i\,(\beta_0 + \beta_1 x_i) - \exp(\beta_0 + \beta_1 x_i) - \ln(y_i!) \right]
\]

Uma propriedade útil: o estimador de Poisson converge para $\beta$ mesmo que
$Y$ não seja estritamente Poisson, contanto que $E[Y\mid x] = \exp(x'\beta)$
(estimador quasi-MLE; Gourieroux, Monfort e Trognon 1984).

\subsection*{Binomial Negativa}

O modelo Poisson impõe $E[Y\mid x] = \mathrm{Var}[Y\mid x]$ (equidispersão).
Quando a variância supera a média (\emph{superdispersão}), o modelo de
Binomial Negativa (NB2) generaliza:

\[
  \mathrm{Var}[Y \mid x] = E[Y\mid x] + \alpha\,(E[Y\mid x])^2
\]

O parâmetro $\alpha > 0$ captura a superdispersão. Poisson é o caso-limite
$\alpha \to 0$.

\subsection*{Zero-Inflated Negative Binomial (ZINB)}

Muitas aplicações apresentam zeros em excesso além do que qualquer modelo de
contagem simples prevê — processos de dois estágios: primeiro decide-se se há
evento ($D=1$ com prob.\ $\pi_i$), depois conta-se a frequência. O ZINB combina:

\[
  P(Y_i = 0) = (1-\pi_i) + \pi_i\,(1+\alpha\lambda_i)^{-1/\alpha}
  \qquad
  P(Y_i = k) = \pi_i\,P_{\mathrm{NB}}(k \mid \lambda_i,\alpha)\quad (k\geq 1)
\]

% ─────────────────────────────────────────────────────────────────────────────
\section{Verificação por DGP}

O DGP cria contagens aproximadamente Poisson com $\lambda_i = e^{0{,}5 + x_i}$,
$x \sim U(0,1)$. O estimador deve recuperar $\beta_0 \approx 0{,}5$ e $\beta_1 \approx 1{,}0$.

\begin{lstlisting}[caption={Poisson, NegBin e ZINB}]
seed(42)
generate df x   = runiform()                         // x ~ U(0,1)
generate df lam = exp(0.5 + 1.0 * x)
generate df u   = runiform()
generate df y   = floor(lam + u)                     // contagem ≥ 0

let m_pois = poisson(y ~ x, df)
display m_pois

let m_nb   = nbreg(y_nb ~ x, df)
display m_nb

let m_zinb = zinb(y_zi ~ x, df)
display m_zinb
\end{lstlisting}

\subsection*{Poisson}

\begin{lstlisting}[language={}, basicstyle=\ttfamily\small,
  frame=none, numbers=none, backgroundcolor=\color{codbg},
  xleftmargin=0pt, xrightmargin=0pt, breaklines=false]
========================= Poisson Regression Results =========================
Dep. Variable:     y  || Log-Likelihood:   102.0744  AIC:  -200.1488
Method:         IRLS  || Pseudo R-sq:        0.7709  N:         500
------------------------------------------------------------------------------
                   coef    std err        z    P>|z|
const            0.5246     0.0602    8.709    0.000      ***
x                0.9733     0.0960   10.141    0.000      ***
\end{lstlisting}

$\hat\beta_0 = 0{,}52 \approx 0{,}5$ e $\hat\beta_1 = 0{,}97 \approx 1{,}0$ —
estimativas próximas dos valores verdadeiros.

\subsection*{Binomial Negativa}

\begin{lstlisting}[language={}, basicstyle=\ttfamily\small,
  frame=none, numbers=none, backgroundcolor=\color{codbg},
  xleftmargin=0pt, xrightmargin=0pt, breaklines=false]
================ Negative Binomial Regression (alpha=0.0010) =================
                   coef    std err        z    P>|z|
const            0.4613     0.0177   26.025    0.000      ***
x                1.1950     0.0276   43.249    0.000      ***
\end{lstlisting}

$\hat\alpha \approx 0{,}001$ — mínima superdispersão, NegBin colapsa para Poisson. Com dados genuinamente superdispersos $\hat\alpha$ é significativamente positivo.

\subsection*{ZINB}

\begin{lstlisting}[language={}, basicstyle=\ttfamily\small,
  frame=none, numbers=none, backgroundcolor=\color{codbg},
  xleftmargin=0pt, xrightmargin=0pt, breaklines=false]
=============== Zero-Inflated Negative Binomial (ZINB) Results ===============
---- Count Model ----
const  0.2994 (0.1347)  **
x      1.2435 (0.2051)  ***
---- Inflate Model (Logit) ----
z0     0.4392 (0.2264)  *
z1     0.5121 (0.3828)
\end{lstlisting}

A equação de inflação de zeros estima a probabilidade de um zero estrutural;
a equação de contagem modela o processo condicional ao evento ocorrer.

% ─────────────────────────────────────────────────────────────────────────────
\section{Sintaxe}

\begin{lstlisting}
poisson(y ~ x1 + x2, df)
nbreg(y ~ x1 + x2, df)
zinb(y ~ x1 + x2, df)
zinb(y ~ x1 + x2, df, inflate = z1 + z2)  // equação de inflação separada
\end{lstlisting}

\begin{center}
\begin{tabular}{ll}
\toprule
\textbf{Estimador} & \textbf{Quando usar} \\
\midrule
\texttt{poisson}  & contagens com equidispersão; quasi-MLE robusto \\
\texttt{nbreg}    & superdispersão ($\hat\alpha$ sig.\ positivo) \\
\texttt{zinb}     & excesso de zeros além do esperado por NB \\
\bottomrule
\end{tabular}
\end{center}

\begin{nota}
  Coeficientes de Poisson e NB são log-razões: $\beta_1 = 1$ implica que um
  aumento unitário em $x$ multiplica a contagem esperada por $e^1 \approx 2{,}72$.
  Para obter \emph{efeitos marginais médios} (AME): use
  \lstinline{margins(modelo)}.
\end{nota}
